\begin{abstract}
智能体强化学习（RL）已经成为云集群中的一种关键新型工作负载，它使大语言模型（LLM）能够通过与真实世界交互来解决复杂问题。然而，与传统 RL 不同，智能体 RL 需要大量位于主训练集群之外的外部云资源，例如用于代码执行的 CPU、用于奖励模型的 GPU 等。现有智能体 RL 框架通常依赖静态过量预留，也就是把资源长期绑定在寿命较长的轨迹上，或者按任务隔离部署，从而带来严重的资源低效。

我们提出动作级编排思想，并将其实现为 \sysname，一个统一的资源管理系统，用于支持细粒度的外部资源共享与弹性调度。\sysname 采用统一的动作级建模和弹性调度算法，在满足异构资源约束的同时最小化动作完成时间（ACT）。进一步地，我们还为不同类型、不同拓扑特征的资源设计了异构资源管理器，以高效支持动作级执行。真实世界智能体 RL 任务上的评测表明，\sysname 可将平均 ACT 最多降低 4.3$\times$，将 RL 训练 step 时长最多加速 1.5$\times$，并最多节省 71.2$\%$ 的外部资源。该系统已经部署并用于支持 MiMo 系列模型的训练。
\end{abstract}
